\begin{figure}[t]
\begin{minted}[fontsize=\small,frame=single,framesep=2mm]{text}
# Teleportation in Stim
# Qubit 0: message qubit (initialized to |0⟩)
# Qubit 1: Alice's EPR qubit
# Qubit 2: Bob's EPR qubit

# Create EPR pair between Alice (q1) and Bob (q2)
H 1
CNOT 1 2

# Alice entangles message qubit with her EPR qubit
CNOT 0 1
H 0

# Alice measures both qubits
M 0 1

# Bob applies corrections based on measurement results
# rec[-1] is measurement of qubit 1 (most recent)
# rec[-2] is measurement of qubit 0 (second most recent)
CX rec[-1] 2
CZ rec[-2] 2

# Result: qubit 2 now holds the original state of qubit 0
\end{minted}
\caption{Teleportation implemented in Stim's assembly-style circuit language. Measurement results are referenced using \texttt{rec[-k]} indices, where \texttt{rec[-1]} refers to the most recent measurement. Conditional Pauli corrections (\texttt{CX rec[-1] 2}) are applied only when the referenced measurement outcome is 1, implementing real-time classical feedback within the stabilizer formalism.}
\label{fig:stim-teleport}
\end{figure}
